\chapter{Genus of a smooth proper curve}
\label{chap:Genus}

Throughout this chapter, \(k\) denotes a field and \(C\) a smooth proper geometrically irreducible curve over \(k\) of relative dimension \(1\). In other words, \(C \to \Spec k\) is a morphism of schemes carrying the geometric typeclasses that make it a curve in the modern sense.

\section{Definition}

% NOTE: classical definition \(g(C) = \dim_k H^1(C, \mathcal O_C)\); see
% Hartshorne, *Algebraic Geometry*, Ch.~IV \S1 (genus of a curve) — TOC
% summarised in `references/hartshorne-ag.md`. The finiteness backing
% (Serre on coherent cohomology of a proper \(k\)-scheme) is queued in
% `references/stacks-0BUG.md` and chapter
% \texttt{Cohomology\_StructureSheafModuleK}.
\begin{definition}[Genus of a smooth proper curve]\leanok
  \label{def:genus}
  \dcref{custom}
  \uses{def:Scheme_HModule, def:Scheme_toModuleKSheaf}
  \lean{AlgebraicGeometry.genus}

  The \emph{genus} of a smooth proper curve \(C\) over \(k\) is the natural number
  \[
    g(C) \;=\; \dim_k H^1(C,\, \mathcal O_C),
  \]
  the \(k\)-dimension of the first sheaf-cohomology group of the structure sheaf of \(C\).
\end{definition}

\begin{proof}

  \leanok
  The definition is realised by passing the structure sheaf through the sheaf-of-\(k\)-modules construction and applying the \(k\)-linear sheaf cohomology \(H^1\). The result is a \(k\)-vector space; its dimension is extracted via the standard finiteness rank.

  \begin{remark}[Lean implementation]
    The formal definition uses the project's \(\ModuleCat k\)-flavoured sheaf cohomology of the structure sheaf:
    \[
      \mathtt{genus}\,C \;:=\; \Module.\mathtt{finrank}\,k\,\bigl(\Scheme.\HModule\,k\,(\Scheme.\toModuleKSheaf\,C)\,1\bigr).
    \]
    The carrier \(\toModuleKSheaf C\) is the sheaf-of-\(k\)-modules constructed from \(\struct C\) via the structure morphism; the wrapper \(\HModule k F 1\) is the \(\Module k\)-flavoured sheaf cohomology, given by \(\Ext\) from the constant sheaf at \(\ModuleCat.\mathrm{of}\,k\,k\), which carries a canonical \(\Module k\)-structure inherited from the \(\Linear k\)-enrichment of the sheaf category.
  \end{remark}

  Mathlib's \(\Module.\mathtt{finrank}\) returns the \(k\)-dimension when the underlying module is finite-dimensional and zero otherwise. For a smooth proper geometrically irreducible curve, Serre's theorem on coherent cohomology of a proper \(k\)-scheme guarantees \(\dim_k H^1(C, \struct C) < \infty\), so \(g(C)\) is the geometric genus on the relevant case. The closure is therefore the one-liner above, confirmed to typecheck against current Mathlib.
\end{proof}

Equivalent reformulations available once Riemann--Roch and Serre duality are formalised:
\begin{itemize}
  \item \(g(C) = \dim_k H^0(C,\, \Omega^1_{C/k})\) (Serre duality).
  \item \(g(C) = 1 - \chi(\mathcal O_C)\), where \(\chi\) is the Euler characteristic.
  \item For a smooth projective curve over \(\mathbb C\), \(g(C)\) equals the topological genus of \(C(\mathbb C)\).
\end{itemize}
Each of these requires Mathlib infrastructure that is currently absent. The first reformulation in particular requires the sheaf of relative differentials of schemes (Phase~B), and the others require Riemann--Roch / Serre duality / a topological-vs-algebraic geometry comparison.

\section{Mathlib gap}
\label{sec:Genus_status}

The honest definition \(\dim_k H^1(C, \mathcal O_C)\) requires a \(\Module k\)-flavoured sheaf cohomology of \(\struct C\) on the topology of opens of the underlying space of \(C\). Through earlier iterations the project has built this:

\begin{itemize}
  \item \textbf{Sheaf cohomology of \(\AddCommGrpCat\)-valued sheaves} (Mathlib): defined as \(\Ext\) from the constant sheaf at \(\AddCommGrpCat.\mathrm{of}\,(\ULift\,\mathbb Z)\).
  \item \textbf{Sheaf cohomology of \(\ModuleCat k\)-valued sheaves} (this project): the parallel \(\Module k\)-valued construction with unit object \(\ModuleCat.\mathrm{of}\,k\,k\) instead of \(\ULift\,\mathbb Z\), gluing through the \(\Linear k\)-enrichment of the sheaf category to recover a \(\Module k\) structure on the cohomology groups.
  \item \textbf{Structure sheaf as a \(\ModuleCat k\)-valued sheaf} (this project): for \(C\) over \(\Spec k\), the sheafification of the presheaf \(U \mapsto \ModuleCat.\mathrm{of}\,k\,(\Gamma(C, U))\) where the \(k\)-action is the algebra structure induced by the structure morphism.
  \item \textbf{\(H^0\) algebraic bridge} (this project): \(H^0(C, \struct C)_{\Module k} \simeq_k \Hom_{\Sheaf}(\mathrm{const}, \toModuleKSheaf C)\), via Mathlib's \(\Ext\) linear equivalence. Used downstream for the geometric content \(H^0(C, \struct C) \simeq k\) on a connected proper \(k\)-curve.
\end{itemize}

The remaining Mathlib gap for the \emph{theorem} (as opposed to the \emph{definition}) of genus:

\begin{itemize}
  \item \textbf{Serre finiteness for \(H^i(C, \mathcal F)\) on a proper \(k\)-curve.} Mathlib's \(\ModuleCat.\mathtt{finite\_ext}\) provides \(\Module.\mathtt{Finite}\) for module-Ext only, not sheaf-Ext. The Čech-cohomology scaffold is the recommended direction.
\end{itemize}

\section{User authorisation of \texttt{noncomputable}}
\label{sec:Genus_noncomputable}

The original challenge file writes \(\mathtt{def\ genus}\), not \(\mathtt{noncomputable\ def\ genus}\). Because \(\Module.\mathtt{finrank}\), \(\Abelian.\Ext.\mathtt{instModule}\), and the structure-sheaf-as-\(k\)-module construction are all noncomputable in Mathlib, the honest closure of the body requires the \(\mathtt{noncomputable}\) modifier on the \(\mathtt{def}\). The user explicitly authorised agents to add this modifier as a non-signature change: names, argument types, and argument order remain verbatim.

\section{What this iteration does and does not}

This iteration \emph{defines} the genus. It does not \emph{compute} it on any specific curve, and it does not prove that the result equals the geometric genus in the classical sense (which requires Serre finiteness for non-vacuity and Riemann--Roch for verification on examples). The downstream theorem \(\mathtt{smoothOfRelativeDimension\_genus}\) in \(\mathtt{Jacobian.lean}\) states that the Jacobian of \(C\) is smooth of relative dimension \(g(C)\) --- that theorem will require Serre finiteness on the way and is queued for the iteration that closes the Jacobian construction itself.
